\documentclass[aspectratio=169]{beamer}

% Tema UFS --- Universidade Federal de Sergipe
\usetheme{UFS}

% Pacotes base
\usepackage[utf8]{inputenc}
\usepackage[portuguese]{babel}
\usepackage{amsmath,amssymb}
\usepackage{graphicx}
\usepackage{booktabs}
\usepackage{xcolor}

% TikZ e bibliotecas
\usepackage{tikz}
\usetikzlibrary{shapes.geometric, arrows.meta, positioning, matrix, fit,
  backgrounds, decorations.pathreplacing, shadows, patterns, calc}

% Listagens --- paleta VS Code Light+
\usepackage{listings}
\definecolor{bgcolor}{HTML}{FFFFFF}
\definecolor{linecolor}{HTML}{DDDDDD}
\definecolor{numcolor}{HTML}{999999}
\definecolor{kwcolor}{HTML}{0000FF}
\definecolor{strcolor}{HTML}{A31515}
\definecolor{cmtcolor}{HTML}{008000}

\lstdefinestyle{ufsStyle}{
  backgroundcolor=\color{bgcolor},
  basicstyle=\ttfamily\footnotesize\color{black},
  keywordstyle=\color{kwcolor}\bfseries,
  stringstyle=\color{strcolor},
  commentstyle=\color{cmtcolor}\itshape,
  numberstyle=\tiny\color{numcolor},
  numbers=left, numbersep=12pt,
  xleftmargin=20pt, framexleftmargin=20pt,
  breaklines=true, tabsize=4,
  keepspaces=true, showspaces=false,
  showstringspaces=false, showtabs=false,
  frame=single, rulecolor=\color{linecolor},
  framesep=4pt, captionpos=b, upquote=true,
  columns=fullflexible,
}
\lstset{style=ufsStyle, language=C}

% Estilo para linha de comando / Makefile (sem numeracao)
\lstdefinestyle{shellStyle}{
  style=ufsStyle, language=bash, numbers=none,
  basicstyle=\ttfamily\scriptsize\color{black},
}
\lstdefinestyle{makeStyle}{
  style=ufsStyle, language=make, numbers=none,
  basicstyle=\ttfamily\scriptsize\color{black},
}

% --- Estilos TikZ reutilizados ---
\tikzset{
  mod/.style={rectangle, draw=UFSBlue, fill=UFSBlue!8, rounded corners,
    minimum width=2.6cm, minimum height=0.85cm, font=\small\ttfamily,
    align=center},
  modh/.style={rectangle, draw=UFSYellow!80!black, fill=UFSYellow!20,
    rounded corners, minimum width=2.2cm, minimum height=0.7cm,
    font=\small\ttfamily, align=center},
  modr/.style={rectangle, draw=UFSRed, fill=UFSRed!8, rounded corners,
    minimum width=2.6cm, minimum height=0.85cm, font=\small\ttfamily,
    align=center},
  obj/.style={rectangle, draw=UFSGray, fill=black!5, minimum width=1.6cm,
    minimum height=0.7cm, font=\small\ttfamily},
  cel/.style={rectangle, draw=UFSBlue, fill=UFSBlue!8, minimum width=1.1cm,
    minimum height=0.75cm, font=\small\ttfamily},
  rot/.style={font=\scriptsize\ttfamily, text=UFSGray},
  leg/.style={font=\scriptsize, text=UFSGray},
  seta/.style={->, thick, UFSBlue},
  setar/.style={->, thick, UFSRed},
  setag/.style={->, thick, UFSGray},
  caixa/.style={rectangle, draw=UFSGray, rounded corners, fill=black!3,
    font=\scriptsize, align=center, inner sep=4pt},
}

% --- Metadados ---
\title{Modularização de programas}
\subtitle{Headers, múltiplos arquivos e bibliotecas}
\author{Prof.\ André Yoshiaki Kashiwabara}
\institute{DCOMP -- Universidade Federal de Sergipe\\
\texttt{andreyoshiaki@dcomp.ufs.br}}
\date{}

\begin{document}

\begin{frame}[plain]
  \titlepage
\end{frame}

\begin{frame}{O que você vai aprender hoje}
  \begin{center}
  \begin{tikzpicture}[
    item/.style={rectangle, draw=UFSBlue, fill=UFSBlue!10, rounded corners,
      minimum width=10cm, minimum height=0.72cm, font=\small}
  ]
    \node[item] (t1) at (0, 0)    {Separar \emph{o que} um módulo oferece do \emph{como} ele faz};
    \node[item] (t2) at (0,-0.95) {O header (\texttt{.h}) como contrato --- e as guardas de inclusão};
    \node[item] (t3) at (0,-1.90) {De \texttt{.c} a executável: pré-processador, compilador, ligador};
    \node[item] (t4) at (0,-2.85) {Um \texttt{Makefile} que recompila só o que mudou};
    \node[item] (t5) at (0,-3.80) {Bibliotecas: estática (\texttt{.a}) e compartilhada (\texttt{.so})};
    \draw[seta] (t1) -- (t2);
    \draw[seta] (t2) -- (t3);
    \draw[seta] (t3) -- (t4);
    \draw[seta] (t4) -- (t5);
  \end{tikzpicture}
  \end{center}
\end{frame}

\begin{frame}{O problema: o arquivo que só cresce}
  \begin{columns}[T]
    \begin{column}{0.5\textwidth}
      \textbf{Onde estamos}\\[0.3em]
      O \texttt{notas.c} da aula passada tem 68 linhas e faz três coisas
      diferentes: cuida da matriz, calcula estatísticas e imprime o relatório.

      \medskip
      Some mais duas funcionalidades e ele passa de 300.
    \end{column}
    \begin{column}{0.5\textwidth}
      \textbf{O que dói quando ele cresce}\\[0.3em]
      \textcolor{UFSRed}{Reuso:} para usar \texttt{cria\_matriz} em outro
      programa, você copia e cola.

      \medskip
      \textcolor{UFSRed}{Recompilação:} mudou uma linha, recompila tudo.

      \medskip
      \textcolor{UFSRed}{Trabalho em equipe:} dois colegas, um arquivo.
    \end{column}
  \end{columns}
  \begin{block}{A pergunta de hoje}
    Como quebrar um programa em pedaços que compilam separado --- e continuam
    formando um executável só?
  \end{block}
\end{frame}

% ==================================================================
\section{O programa de hoje}
% ==================================================================

\begin{frame}[fragile]{O \texttt{notas.c} de hoje, ainda em um arquivo só}
\begin{lstlisting}[basicstyle=\ttfamily\scriptsize]
#include <stdio.h>
#include <stdlib.h>
#include <math.h>

double **cria_matriz(int nl, int nc)     { /* ... */ }
void libera_matriz(double **m, int nl)   { /* ... */ }

double media(const double *v, int n)         { /* ... */ }
double desvio_padrao(const double *v, int n) { /* ... */ }

int main(void) { /* monta a turma, imprime media e desvio */ }
\end{lstlisting}
  \begin{block}{}
    \small Três assuntos convivendo: \textbf{memória}, \textbf{matemática} e
    \textbf{apresentação}. Cada um muda por um motivo diferente.
  \end{block}
\end{frame}

\begin{frame}{A quebra: quatro arquivos, três responsabilidades}
  \begin{center}
  \begin{tikzpicture}[node distance=0.6cm]
    \node[modr] (main) at (0,0) {main.c\\[-2pt]\scriptsize monta e imprime};
    \node[mod] (turma) at (0,-1.6) {turma.c\\[-2pt]\scriptsize a turma};
    \node[mod] (matriz) at (-3.2,-3.2) {matriz.c\\[-2pt]\scriptsize memória};
    \node[mod] (estat) at (3.2,-3.2) {estat.c\\[-2pt]\scriptsize matemática};
    \draw[seta] (main) -- (turma);
    \draw[seta] (turma) -- (matriz);
    \draw[seta] (turma) -- (estat);
    \node[leg, anchor=west] at (1.6,-0.8) {``usa''};
  \end{tikzpicture}
  \end{center}
  \begin{block}{}
    \small A seta vai de quem \emph{usa} para quem \emph{oferece}. Repare que
    \texttt{matriz.c} e \texttt{estat.c} não sabem que a turma existe --- por
    isso servem para qualquer outro programa.
  \end{block}
\end{frame}

\begin{frame}[fragile]{\texttt{matriz.h} e \texttt{matriz.c}}
  \begin{columns}[T]
    \begin{column}{0.46\textwidth}
      \textbf{\texttt{matriz.h}} --- o que existe
\begin{lstlisting}[basicstyle=\ttfamily\tiny, numbers=none]
#ifndef MATRIZ_H
#define MATRIZ_H

double **cria_matriz(int nl, int nc);
void libera_matriz(double **m, int nl);

#endif /* MATRIZ_H */
\end{lstlisting}
    \end{column}
    \begin{column}{0.52\textwidth}
      \textbf{\texttt{matriz.c}} --- como funciona
\begin{lstlisting}[basicstyle=\ttfamily\tiny, numbers=none]
#include <stdlib.h>
#include "matriz.h"

double **cria_matriz(int nl, int nc)
{
    double **m = malloc(nl * sizeof(double *));
    /* ... o mesmo codigo da aula passada ... */
    return m;
}
\end{lstlisting}
    \end{column}
  \end{columns}
  \begin{block}{}
    \small Sete linhas contra sessenta: o header cabe na tela, e é ele que o
    resto do programa lê.
  \end{block}
\end{frame}

\begin{frame}[fragile]{\texttt{estat.h} e \texttt{estat.c}}
  \begin{columns}[T]
    \begin{column}{0.46\textwidth}
      \textbf{\texttt{estat.h}}
\begin{lstlisting}[basicstyle=\ttfamily\tiny, numbers=none]
#ifndef ESTAT_H
#define ESTAT_H

double media(const double *v, int n);
double desvio_padrao(const double *v, int n);

#endif /* ESTAT_H */
\end{lstlisting}
    \end{column}
    \begin{column}{0.52\textwidth}
      \textbf{\texttt{estat.c}}
\begin{lstlisting}[basicstyle=\ttfamily\tiny, numbers=none]
#include <math.h>        /* sqrt */
#include "estat.h"

double desvio_padrao(const double *v, int n)
{
    double mu = media(v, n), s = 0.0;
    /* ... */
    return n > 0 ? sqrt(s / n) : 0.0;
}
\end{lstlisting}
    \end{column}
  \end{columns}
  \begin{block}{}
    \small Guarde o \texttt{sqrt}: ele vai reaparecer no fim da aula, quando
    falarmos da biblioteca padrão.
  \end{block}
\end{frame}

\begin{frame}[fragile]{\texttt{turma.h}: o header que também define um tipo}
\begin{lstlisting}[basicstyle=\ttfamily\scriptsize]
#ifndef TURMA_H
#define TURMA_H

typedef struct {
    double **notas;
    int nalunos;
    int navaliacoes;
} Turma;

Turma *turma_cria(int nalunos, int navaliacoes);
void turma_libera(Turma *t);
void turma_relatorio(const Turma *t);

#endif /* TURMA_H */
\end{lstlisting}
\end{frame}

\begin{frame}[fragile]{\texttt{turma.c}: quem usa os outros dois}
\begin{lstlisting}[basicstyle=\ttfamily\scriptsize]
#include <stdio.h>
#include "turma.h"
#include "matriz.h"
#include "estat.h"
static void cabecalho(void)          /* so turma.c enxerga */
{
    printf("%-8s %8s %8s\n", "aluno", "media", "desvio");
}

void turma_relatorio(const Turma *t)
{
    cabecalho();
    for (int i = 0; i < t->nalunos; i++)
        printf("%-8d %8.2f %8.2f\n", i, media(t->notas[i], t->navaliacoes),
               desvio_padrao(t->notas[i], t->navaliacoes));
}
\end{lstlisting}
\end{frame}

\begin{frame}[fragile]{\texttt{main.c}: um \texttt{include} só}
\begin{lstlisting}[basicstyle=\ttfamily\scriptsize]
#include <stdio.h>
#include "turma.h"       /* nao inclui matriz.h nem estat.h */
int main(void)
{
    double valores[3][4] = { { 8.0, 6.5, 9.0, 7.0 },
                             { 5.0, 7.0, 4.0, 7.0 },
                             { 9.0, 10.0, 8.5, 9.5 } };
    Turma *t = turma_cria(3, 4);
    if (t == NULL) { fprintf(stderr, "sem memoria\n"); return 1; }
    for (int i = 0; i < t->nalunos; i++)
        for (int j = 0; j < t->navaliacoes; j++)
            t->notas[i][j] = valores[i][j];
    turma_relatorio(t);
    turma_libera(t);
    return 0;
}
\end{lstlisting}
\end{frame}

\begin{frame}{O que vai acontecer?}
  \begin{block}{Antes de compilar, escreva sua resposta}
    \begin{enumerate}
      \item \texttt{gcc -c matriz.c} produz o quê? Esse arquivo roda?
      \item \texttt{main.c} não inclui \texttt{matriz.h}. Ainda assim o
        executável final vai conter \texttt{cria\_matriz}?
      \item O que acontece se você compilar \texttt{gcc matriz.c estat.c
        turma.c -o notas}, esquecendo \texttt{main.c}?
      \item E se esquecer o \texttt{-lm} na hora de juntar tudo?
    \end{enumerate}
  \end{block}
\end{frame}

\begin{frame}[fragile]{Compile em duas etapas e compare}
\begin{lstlisting}[style=shellStyle]
$ gcc -Wall -Wextra -c matriz.c estat.c turma.c main.c
$ ls *.o
estat.o  main.o  matriz.o  turma.o
$ gcc matriz.o estat.o turma.o main.o -o notas -lm
$ ./notas
aluno       media   desvio
0            7.62     0.96
1            5.75     1.30
2            9.25     0.56
\end{lstlisting}
  \begin{block}{}
    \small A saída é a mesma da aula passada. O que mudou não foi o programa
    --- foi o \emph{caminho} até o executável.
  \end{block}
\end{frame}

% ==================================================================
\section{Interface e implementação}
% ==================================================================

\begin{frame}{Declarar não é definir}
  \begin{block}{Definição}
    \textbf{Declaração} anuncia que um nome existe e qual é o seu tipo:
    \texttt{double media(const double *v, int n);}

    \smallskip
    \textbf{Definição} fornece o corpo --- o código que será executado.
  \end{block}
  \begin{exampleblock}{Intuição}
    A declaração é a ficha do cardápio: nome do prato, o que entra, o que sai.
    A definição é a cozinha. Quem faz o pedido só precisa do cardápio.
  \end{exampleblock}
\end{frame}

\begin{frame}{O header é um contrato entre dois lados}
  \begin{center}
  \begin{tikzpicture}
    \node[modr] (usa) at (-4,0) {turma.c\\[-2pt]\scriptsize chama};
    \node[modh] (h) at (0,0) {matriz.h\\[-2pt]\scriptsize o contrato};
    \node[mod] (impl) at (4,0) {matriz.c\\[-2pt]\scriptsize cumpre};
    \draw[seta] (usa) -- (h) node[midway, above, leg] {confia};
    \draw[seta] (impl) -- (h) node[midway, above, leg] {obedece};
    \node[caixa, text width=10cm] at (0,-1.7) {%
      Os dois lados leem o \emph{mesmo} arquivo. Por isso o compilador
      consegue avisar quando a chamada e a implementação discordam --- desde
      que os dois incluam o header.};
  \end{tikzpicture}
  \end{center}
\end{frame}

\begin{frame}{O que vai no \texttt{.h} e o que fica no \texttt{.c}}
  \begin{center}
  \small
  \begin{tabular}{@{}ll@{}}
    \toprule
    \textbf{No header (\texttt{.h})} & \textbf{No fonte (\texttt{.c})} \\
    \midrule
    protótipos de funções públicas   & corpo das funções \\
    \texttt{typedef} e \texttt{struct} usados na interface & funções auxiliares \texttt{static} \\
    \texttt{enum} e \texttt{\#define} de constantes & variáveis do módulo (\texttt{static}) \\
    \texttt{extern} de variáveis globais & a \emph{definição} dessas globais \\
    \bottomrule
  \end{tabular}
  \end{center}
  \begin{block}{Regra prática}
    \small Se remover a linha do header quebra quem usa o módulo, ela pertence
    ao header. Se não quebra, ela é detalhe de implementação.
  \end{block}
\end{frame}

\begin{frame}[fragile]{\texttt{\#include} é copiar e colar}
\begin{lstlisting}[style=shellStyle]
$ cat mini.c
#include "matriz.h"
int main(void) { return 0; }

$ gcc -E mini.c            # so o pre-processador
# 1 "mini.c"
# 1 "./matriz.h" 1
double **cria_matriz(int nl, int nc);
void libera_matriz(double **m, int nl);
# 2 "mini.c" 2
int main(void) { return 0; }
\end{lstlisting}
  \begin{block}{}
    \small O pré-processador não entende C: ele troca a linha
    \texttt{\#include} pelo conteúdo do arquivo. As linhas com \texttt{\#}
    apenas registram de onde veio cada trecho. Com \texttt{<stdio.h>} no lugar,
    a saída passa de 500 linhas.
  \end{block}
\end{frame}

\begin{frame}{\texttt{"aspas"} ou \texttt{<}ângulos\texttt{>}?}
  \begin{columns}[T]
    \begin{column}{0.5\textwidth}
      \textbf{\texttt{\#include "matriz.h"}}\\[0.3em]
      Procura primeiro no diretório do arquivo que incluiu; depois nos
      diretórios do sistema.

      \medskip
      Use para os \emph{seus} headers.
    \end{column}
    \begin{column}{0.5\textwidth}
      \textbf{\texttt{\#include <stdio.h>}}\\[0.3em]
      Procura só nos diretórios do sistema (e nos que você acrescentar com
      \texttt{-I}).

      \medskip
      Use para a biblioteca padrão e para bibliotecas instaladas.
    \end{column}
  \end{columns}
  \begin{block}{}
    \small \texttt{gcc -Icaminho/include} acrescenta um diretório à busca ---
    é assim que você usa headers que não estão ao lado do seu \texttt{.c}.
  \end{block}
\end{frame}

\begin{frame}{O problema: o mesmo header, duas vezes}
  \begin{center}
  \begin{tikzpicture}
    \node[modr] (c) at (0,0) {turma.c};
    \node[modh] (t) at (-2.6,-1.5) {turma.h};
    \node[modh] (e) at (2.6,-1.5) {estat.h};
    \node[modh] (m) at (0,-3.0) {tipos.h};
    \draw[seta] (c) -- (t);
    \draw[seta] (c) -- (e);
    \draw[seta] (t) -- (m);
    \draw[seta] (e) -- (m);
    \node[caixa, text width=5.2cm, anchor=west] at (3.6,-2.0) {%
      \texttt{tipos.h} entra \textbf{duas vezes} no texto que o compilador
      recebe.\\[3pt]
      Se ele definir um \texttt{typedef} ou uma \texttt{struct}:
      \textcolor{UFSRed}{\texttt{redefinition of `Ponto'}}};
  \end{tikzpicture}
  \end{center}
\end{frame}

\begin{frame}[fragile]{A solução: guardas de inclusão}
  \begin{columns}[T]
    \begin{column}{0.54\textwidth}
\begin{lstlisting}[basicstyle=\ttfamily\scriptsize, numbers=none]
#ifndef MATRIZ_H     /* se ainda nao definido */
#define MATRIZ_H     /* marque como definido  */

double **cria_matriz(int nl, int nc);
void libera_matriz(double **m, int nl);

#endif /* MATRIZ_H */
\end{lstlisting}
    \end{column}
    \begin{column}{0.44\textwidth}
      \small
      Na segunda inclusão, \texttt{MATRIZ\_H} já existe: tudo entre
      \texttt{\#ifndef} e \texttt{\#endif} é descartado.

      \medskip
      O nome da macro precisa ser único no projeto --- costume:
      \texttt{NOMEDOARQUIVO\_H}.
    \end{column}
  \end{columns}
  \begin{block}{}
    \small \texttt{\#pragma once} faz o mesmo em uma linha e é aceito por
    todos os compiladores atuais, mas não é parte do padrão C. Prefira as
    guardas quando o código precisar ser portável.
  \end{block}
\end{frame}

\begin{frame}{Recapitulando: interface e implementação}
  \begin{block}{O que aprendemos}
    \begin{itemize}
      \item declarar anuncia; definir fornece o corpo --- o header carrega as
        declarações;
      \item quem usa e quem implementa leem o mesmo \texttt{.h}: é isso que
        permite ao compilador conferir as chamadas;
      \item \texttt{\#include} é substituição de texto, feita antes de o
        compilador ver qualquer C;
      \item toda \texttt{.h} sua começa com \texttt{\#ifndef}/\texttt{\#define}
        e termina com \texttt{\#endif}.
    \end{itemize}
  \end{block}
\end{frame}

% ==================================================================
\section{De .c a executável}
% ==================================================================

\begin{frame}{As quatro etapas escondidas atrás de \texttt{gcc}}
  \begin{center}
  \begin{tikzpicture}[node distance=0.35cm]
    \node[obj, minimum width=1.5cm] (c) at (0,0) {.c};
    \node[obj, minimum width=1.5cm, right=1.1cm of c] (i) {.i};
    \node[obj, minimum width=1.5cm, right=1.1cm of i] (s) {.s};
    \node[obj, minimum width=1.5cm, right=1.1cm of s] (o) {.o};
    \node[obj, minimum width=1.7cm, right=1.1cm of o, draw=UFSRed, fill=UFSRed!8] (x) {notas};
    \draw[seta] (c) -- (i) node[midway, above, leg] {\texttt{cpp}};
    \draw[seta] (i) -- (s) node[midway, above, leg] {compila};
    \draw[seta] (s) -- (o) node[midway, above, leg] {monta};
    \draw[seta] (o) -- (x) node[midway, above, leg] {liga};
    \node[leg] at (1.75,-0.75) {\texttt{-E}};
    \node[leg] at (4.05,-0.75) {\texttt{-S}};
    \node[leg] at (6.35,-0.75) {\texttt{-c}};
    \node[leg] at (8.75,-0.75) {\texttt{ld}};
  \end{tikzpicture}
  \end{center}
  \begin{block}{}
    \small Cada flag \emph{para} o processo em um ponto. \texttt{gcc arquivo.c}
    sem flag nenhuma percorre as quatro etapas de uma vez --- por isso parece
    um passo só.
  \end{block}
\end{frame}

\begin{frame}{Unidade de tradução: o compilador vê um arquivo por vez}
  \begin{columns}[T]
    \begin{column}{0.52\textwidth}
      \textbf{O que o compilador enxerga ao ler \texttt{turma.c}}\\[0.3em]
      \small O texto de \texttt{turma.c} depois dos \texttt{include}s --- e
      mais nada. Ele \emph{não} abre \texttt{matriz.c}.
    \end{column}
    \begin{column}{0.46\textwidth}
      \textbf{Consequência}\\[0.3em]
      \small Sobre \texttt{cria\_matriz} ele sabe só o que o header declarou:
      o tipo. Basta para gerar a chamada; não basta para saber onde o código
      está.
    \end{column}
  \end{columns}
  \begin{block}{Quem resolve o resto}
    O \textbf{ligador}, depois, com todos os \texttt{.o} na mão. É por isso
    que erros de nome errado aparecem tão tarde.
  \end{block}
\end{frame}

\begin{frame}[fragile]{O que sobra dentro de um \texttt{.o}}
\begin{lstlisting}[style=shellStyle]
$ nm turma.o
0000000000000000 T turma_cria          # T = definido aqui (publico)
0000000000000090 T turma_libera
00000000000000d0 T turma_relatorio
0000000000000160 t cabecalho           # t = definido, mas so daqui
                 U cria_matriz         # U = usado, definido em outro lugar
                 U desvio_padrao
                 U media
                 U printf
\end{lstlisting}
  \begin{block}{}
    \small Um \texttt{.o} é código de máquina \emph{mais} duas listas: o que
    ele oferece (\texttt{T}) e o que ele ainda deve (\texttt{U}).
  \end{block}
\end{frame}

\begin{frame}{O ligador casa dívidas com ofertas}
  \begin{center}
  \begin{tikzpicture}
    \node[obj, minimum width=2.0cm] (m)   at (-4.2, 1.1) {main.o};
    \node[obj, minimum width=2.0cm] (t)   at (-4.2, 0.0) {turma.o};
    \node[obj, minimum width=2.0cm] (mat) at (-4.2,-1.1) {matriz.o};
    \node[rot, anchor=west] (u1) at (-2.9, 1.1) {U turma\_cria};
    \node[rot, anchor=west] (t1) at (-2.9, 0.0) {T turma\_cria};
    \node[rot, anchor=west] (u2) at (0.3, 0.0) {U cria\_matriz};
    \node[rot, anchor=west] (t2) at (0.3,-1.1) {T cria\_matriz};
    \draw[setar] (u1.east) to[out=0, in=90] ($(t1.east)+(0.6,0)$) to[out=270, in=0] (t1.east);
    \draw[setar] (u2.east) to[out=0, in=90] ($(t2.east)+(0.6,0)$) to[out=270, in=0] (t2.east);
    \node[caixa, text width=4.6cm] at (3.9, 1.1) {%
      Toda linha \texttt{U} precisa encontrar exatamente \textbf{uma} linha
      \texttt{T} entre os arquivos ligados.};
  \end{tikzpicture}
  \end{center}
  \begin{block}{}
    \small Nenhuma linha \texttt{U} sobrando, nenhuma \texttt{T} repetida:
    é esse o fecho de contas que decide se o executável nasce.
  \end{block}
\end{frame}

\begin{frame}[fragile]{Erro 1: \texttt{undefined reference} --- faltou a oferta}
\begin{lstlisting}[style=shellStyle]
$ gcc main.o -o notas
/usr/bin/ld: main.o: in function `main':
main.c:(.text+0x2a): undefined reference to `turma_cria'
collect2: error: ld returned 1 exit status
\end{lstlisting}
  \begin{block}{As três causas, em ordem de frequência}
    \small
    \begin{enumerate}
      \item esqueceu um \texttt{.o} (ou um \texttt{.c}) na linha de comando;
      \item escreveu o nome diferente no \texttt{.h} e no \texttt{.c}
        (\texttt{turma\_cria} $\times$ \texttt{cria\_turma});
      \item a função existe, mas está em uma biblioteca que você não pediu
        para ligar (\texttt{-lm}, por exemplo).
    \end{enumerate}
  \end{block}
\end{frame}

\begin{frame}[fragile]{Erro 2: \texttt{multiple definition} --- ofertas demais}
  \begin{columns}[T]
    \begin{column}{0.48\textwidth}
      \textbf{O que provoca}
\begin{lstlisting}[basicstyle=\ttfamily\tiny, numbers=none]
/* config.h  -- ERRADO */
int contador = 0;   /* definicao! */
\end{lstlisting}
      \small Incluído em dois \texttt{.c}, vira duas definições da mesma
      variável.
    \end{column}
    \begin{column}{0.50\textwidth}
      \textbf{O jeito certo}
\begin{lstlisting}[basicstyle=\ttfamily\tiny, numbers=none]
/* config.h */
extern int contador;   /* so anuncia */

/* config.c */
int contador = 0;      /* uma definicao */
\end{lstlisting}
    \end{column}
  \end{columns}
\begin{lstlisting}[style=shellStyle]
/usr/bin/ld: turma.o:(.bss+0x0): multiple definition of `contador';
main.o:(.bss+0x0): first defined here
\end{lstlisting}
\end{frame}

\begin{frame}[fragile]{\texttt{static}: o que não sai do arquivo}
\begin{lstlisting}[basicstyle=\ttfamily\scriptsize, numbers=none]
static void cabecalho(void) { ... }   /* funcao so de turma.c   */
static int chamadas = 0;              /* variavel so de turma.c */
\end{lstlisting}
  \begin{columns}[T]
    \begin{column}{0.5\textwidth}
      \textbf{O que muda}\\[0.2em]
      \small O símbolo sai do \texttt{.o} como \texttt{t} minúsculo: o
      ligador não o oferece a ninguém.
    \end{column}
    \begin{column}{0.5\textwidth}
      \textbf{Por que usar}\\[0.2em]
      \small Dois módulos podem ter cada um a sua \texttt{cabecalho} sem
      colidir, e você deixa claro o que é detalhe interno.
    \end{column}
  \end{columns}
  \begin{block}{}
    \small Dentro de uma função, \texttt{static} quer dizer outra coisa: a
    variável sobrevive entre as chamadas. Mesma palavra, dois significados.
  \end{block}
\end{frame}

\begin{frame}{Recapitulando: compilação separada}
  \begin{block}{O que aprendemos}
    \begin{itemize}
      \item \texttt{-c} para em \texttt{.o}; sem flag, \texttt{gcc} vai até o
        executável;
      \item cada \texttt{.c} é compilado sozinho --- o compilador confia no
        header;
      \item o \texttt{.o} lista o que oferece (\texttt{T}) e o que deve
        (\texttt{U}); o ligador faz o casamento;
      \item \texttt{undefined reference} = faltou; \texttt{multiple definition}
        = sobrou; \texttt{static} esconde.
    \end{itemize}
  \end{block}
\end{frame}

% ==================================================================
\section{Automatizar: make}
% ==================================================================

\begin{frame}{Por que não digitar tudo de novo}
  \begin{center}
  \begin{tikzpicture}
    \node[obj, minimum width=1.9cm, draw=UFSRed, fill=UFSRed!8] (ex) at (0,2.1) {notas};
    \node[obj, minimum width=1.6cm] (mo) at (-3.6,0.8) {main.o};
    \node[obj, minimum width=1.6cm] (to) at (0,0.8) {turma.o};
    \node[obj, minimum width=1.6cm] (eo) at (3.6,0.8) {estat.o};
    \node[caixa, text width=2.3cm] (ms) at (-3.6,-0.9)
      {\ttfamily main.c\\ turma.h};
    \node[caixa, text width=2.3cm] (ts) at (0,-0.9)
      {\ttfamily turma.c\\ turma.h\\ matriz.h \ estat.h};
    \node[caixa, text width=2.3cm] (es) at (3.6,-0.9)
      {\ttfamily estat.c\\ estat.h};
    \draw[seta] (ms) -- (mo);
    \draw[seta] (ts) -- (to);
    \draw[seta] (es) -- (eo);
    \draw[setar] (mo) -- (ex);
    \draw[setar] (to) -- (ex);
    \draw[setar] (eo) -- (ex);
  \end{tikzpicture}
  \end{center}
  \begin{block}{}
    \small Mudou \texttt{estat.c}: refazer \texttt{estat.o} e religar. Mudou
    \texttt{turma.h}: \texttt{main.o} e \texttt{turma.o} caem junto. Quem
    decide isso pela data dos arquivos é o \texttt{make}.
  \end{block}
\end{frame}

\begin{frame}[fragile]{Anatomia de uma regra}
\begin{lstlisting}[style=makeStyle]
alvo: pre-requisitos
<TAB>receita

estat.o: estat.c estat.h
	gcc -Wall -Wextra -c estat.c
\end{lstlisting}
  \begin{block}{A regra de ouro do \texttt{make}}
    \small Se o \textbf{alvo} não existe, ou se algum \textbf{pré-requisito} é
    mais recente que ele, execute a \textbf{receita}. Só isso.
  \end{block}
  \begin{alertblock}{}
    \small A receita precisa começar com TAB de verdade --- espaços dão
    \texttt{missing separator}. É o erro número um de quem começa.
  \end{alertblock}
\end{frame}

\begin{frame}[fragile]{O \texttt{Makefile} do \texttt{notas}}
\begin{lstlisting}[style=makeStyle]
CC      = gcc
CFLAGS  = -Wall -Wextra -g
LDLIBS  = -lm
OBJ     = main.o turma.o matriz.o estat.o

notas: $(OBJ)
	$(CC) $(OBJ) -o notas $(LDLIBS)

main.o:   main.c   turma.h
turma.o:  turma.c  turma.h matriz.h estat.h
matriz.o: matriz.c matriz.h
estat.o:  estat.c  estat.h

.PHONY: clean
clean:
	rm -f $(OBJ) notas
\end{lstlisting}
\end{frame}

\begin{frame}[fragile]{As regras sem receita não estão incompletas}
  \begin{columns}[T]
    \begin{column}{0.5\textwidth}
      \textbf{O que o \texttt{make} já sabe}\\[0.3em]
      \small Existe uma regra implícita para \texttt{\%.o: \%.c}, que roda
      \texttt{\$(CC) \$(CFLAGS) -c}.
    \end{column}
    \begin{column}{0.5\textwidth}
      \textbf{O que você acrescenta}\\[0.3em]
      \small Só a informação que o \texttt{make} não tem como adivinhar: de
      quais \emph{headers} cada objeto depende.
    \end{column}
  \end{columns}
  \begin{block}{\texttt{.PHONY}}
    \small \texttt{clean} não é um arquivo. Sem \texttt{.PHONY: clean}, se
    alguém criar um arquivo chamado \texttt{clean}, o \texttt{make} passa a
    dizer ``\texttt{clean is up to date}'' e não limpa mais nada.
  \end{block}
\end{frame}

\begin{frame}[fragile]{\texttt{make} em ação}
\begin{lstlisting}[style=shellStyle]
$ make                          # primeira vez: compila tudo
gcc -Wall -Wextra -g   -c -o main.o main.c
gcc -Wall -Wextra -g   -c -o turma.o turma.c
gcc -Wall -Wextra -g   -c -o matriz.o matriz.c
gcc -Wall -Wextra -g   -c -o estat.o estat.c
gcc main.o turma.o matriz.o estat.o -o notas -lm

$ make                          # nada mudou
make: 'notas' is up to date.

$ touch estat.c && make         # so o que depende de estat.c
gcc -Wall -Wextra -g   -c -o estat.o estat.c
gcc main.o turma.o matriz.o estat.o -o notas -lm
\end{lstlisting}
\end{frame}

\begin{frame}{Recapitulando: make}
  \begin{block}{O que aprendemos}
    \begin{itemize}
      \item o \texttt{make} compara datas: alvo mais velho que
        pré-requisito $\Rightarrow$ refaz;
      \item a receita vai depois de um TAB;
      \item listar os headers de cada objeto é o que faz o
        \texttt{make} acertar quando uma interface muda;
      \item \texttt{.PHONY} marca alvos que não são arquivos.
    \end{itemize}
  \end{block}
\end{frame}

% ==================================================================
\section{Bibliotecas}
% ==================================================================

\begin{frame}{O que é uma biblioteca}
  \begin{block}{Definição}
    Uma \textbf{biblioteca} é um pacote de arquivos-objeto já compilados,
    distribuído junto com os headers que descrevem o que ela oferece.
  \end{block}
  \begin{exampleblock}{Intuição}
    É o passo seguinte da modularização: em vez de entregar quatro
    \texttt{.o} soltos e torcer para ninguém esquecer um, você entrega
    \emph{um} arquivo --- e um \texttt{.h} dizendo o que tem dentro.
  \end{exampleblock}
\end{frame}

\begin{frame}[fragile]{Biblioteca estática: \texttt{ar rcs}}
\begin{lstlisting}[style=shellStyle]
$ gcc -Wall -c matriz.c estat.c turma.c
$ ar rcs libnotas.a matriz.o estat.o turma.o     # r=inserir c=criar s=indice
$ ar t libnotas.a
matriz.o
estat.o
turma.o
$ gcc main.c -I. -L. -lnotas -lm -o notas
\end{lstlisting}
  \begin{block}{Os três flags da ligação}
    \small \texttt{-I.} onde procurar headers \quad
    \texttt{-L.} onde procurar bibliotecas \quad
    \texttt{-lnotas} qual usar --- o nome do arquivo é
    \texttt{lib}\emph{notas}\texttt{.a}.
  \end{block}
\end{frame}

\begin{frame}{O ligador copia só o que você usa}
  \begin{center}
  \begin{tikzpicture}
    \node[caixa, minimum width=2.4cm, minimum height=2.4cm, fill=UFSYellow!15,
      draw=UFSYellow!80!black] (lib) at (-3.6,0) {};
    \node[leg] at (-3.6,1.5) {\texttt{libnotas.a}};
    \node[obj, minimum width=1.5cm] at (-3.6,0.7) {matriz.o};
    \node[obj, minimum width=1.5cm] at (-3.6,0.0) {estat.o};
    \node[obj, minimum width=1.5cm] at (-3.6,-0.7) {turma.o};
    \node[caixa, minimum width=2.6cm, minimum height=2.4cm, draw=UFSRed,
      fill=UFSRed!6] (exe) at (2.6,0) {};
    \node[leg] at (2.6,1.5) {executável \texttt{notas}};
    \node[obj, minimum width=1.5cm] at (2.6,0.7) {main.o};
    \node[obj, minimum width=1.5cm] at (2.6,0.0) {turma.o};
    \node[obj, minimum width=1.5cm] at (2.6,-0.7) {matriz.o};
    \draw[setar] (-2.2,0) -- (1.1,0) node[midway, above, leg] {cópia};
  \end{tikzpicture}
  \end{center}
  \begin{block}{}
    \small Se \texttt{estat.o} não fosse usado por ninguém, ele ficaria de
    fora. O executável final é autossuficiente: a \texttt{.a} pode ser
    apagada depois.
  \end{block}
\end{frame}

\begin{frame}[fragile]{Biblioteca compartilhada: \texttt{-fPIC -shared}}
\begin{lstlisting}[style=shellStyle]
$ gcc -fPIC -c matriz.c estat.c turma.c       # codigo independente de posicao
$ gcc -shared -o libnotas.so matriz.o estat.o turma.o -lm
$ gcc main.c -I. -L. -lnotas -o notas
$ ./notas
./notas: error while loading shared libraries: libnotas.so:
cannot open shared object file: No such file or directory
$ LD_LIBRARY_PATH=. ./notas                   # agora sim
\end{lstlisting}
  \begin{block}{}
    \small Ligou sem reclamar e mesmo assim falhou ao rodar: a
    \texttt{.so} não entra no executável, fica só o \emph{nome} dela. Quem a
    procura, e não acha, é o carregador do sistema.
  \end{block}
\end{frame}

\begin{frame}{Duas horas diferentes de resolver}
  \begin{center}
  \begin{tikzpicture}[scale=0.95, every node/.style={transform shape}]
    \node[leg] at (-3.4,1.6) {\textbf{estática}};
    \node[caixa, minimum width=2.8cm, minimum height=1.5cm, draw=UFSRed,
      fill=UFSRed!6] (e1) at (-3.4,0.4) {main.o\\+ turma.o\\+ matriz.o};
    \node[leg, text width=3cm, align=center] at (-3.4,-1.2)
      {tudo dentro; roda sozinho};
    \node[leg] at (1.4,1.6) {\textbf{compartilhada}};
    \node[caixa, minimum width=2.2cm, minimum height=1.5cm, draw=UFSRed,
      fill=UFSRed!6] (e2) at (0.8,0.4) {main.o\\+ ``preciso de\\libnotas.so''};
    \node[caixa, minimum width=1.9cm, minimum height=1.0cm,
      draw=UFSYellow!80!black, fill=UFSYellow!15] (so) at (4.2,0.4) {libnotas.so};
    \draw[setag, dashed] (e2) -- (so) node[midway, above, leg] {ao rodar};
    \node[leg, text width=4.4cm, align=center] at (2.5,-1.2)
      {o arquivo precisa estar lá na hora da execução};
  \end{tikzpicture}
  \end{center}
\end{frame}

\begin{frame}{Qual das duas usar}
  \begin{center}
  \small
  \begin{tabular}{@{}lll@{}}
    \toprule
     & \textbf{Estática (\texttt{.a})} & \textbf{Compartilhada (\texttt{.so})} \\
    \midrule
    Quando entra   & na ligação        & ao executar \\
    Executável     & maior, completo   & menor, depende do arquivo \\
    Corrigir um bug& religar tudo      & trocar a \texttt{.so} \\
    Dez programas  & dez cópias no disco & uma cópia na memória \\
    Risco          & versão congelada  & sumir/mudar sem avisar \\
    \bottomrule
  \end{tabular}
  \end{center}
  \begin{block}{}
    \small \texttt{ldd notas} lista de quais \texttt{.so} um executável
    depende (no macOS, \texttt{otool -L}).
  \end{block}
\end{frame}

\begin{frame}[fragile]{Você já usa bibliotecas: \texttt{-lm}}
\begin{lstlisting}[style=shellStyle]
$ gcc main.o turma.o matriz.o estat.o -o notas
/usr/bin/ld: estat.o: undefined reference to symbol 'sqrt@@GLIBC_2.2.5'
collect2: error: ld returned 1 exit status
$ gcc main.o turma.o matriz.o estat.o -o notas -lm     # resolvido
\end{lstlisting}
  \begin{block}{\texttt{math.h} e \texttt{libm} são coisas distintas}
    \small O header traz a \emph{declaração} de \texttt{sqrt}; o código está
    em \texttt{libm}. Incluir o header engana o compilador, não o ligador.
    Já \texttt{printf} não pede flag porque \texttt{libc} é ligada sempre.
  \end{block}
  \begin{alertblock}{}
    \small Em alguns sistemas (macOS, glibc recente) a \texttt{libm} já vem
    junto e o erro não aparece. Escreva o \texttt{-lm} assim mesmo.
  \end{alertblock}
\end{frame}

\begin{frame}[fragile]{A ordem dos \texttt{-l} importa}
\begin{lstlisting}[style=shellStyle]
$ gcc -L. -lnotas main.c -o notas       # ERRADO: undefined reference
$ gcc main.c -L. -lnotas -o notas       # certo
\end{lstlisting}
  \begin{block}{Por quê}
    \small O ligador lê a linha de comando da esquerda para a direita
    acumulando dívidas. Quando ele abriu \texttt{libnotas} no primeiro
    comando, ninguém ainda devia \texttt{turma\_cria} --- e ele seguiu adiante
    sem pegar nada.
  \end{block}
  \begin{block}{Regra}
    \small Quem \emph{usa} vem antes de quem \emph{fornece}: primeiro seus
    arquivos, depois as bibliotecas.
  \end{block}
\end{frame}

\begin{frame}{Recapitulando: bibliotecas}
  \begin{block}{O que aprendemos}
    \begin{itemize}
      \item biblioteca = \texttt{.o} empacotados + headers que os descrevem;
      \item estática: \texttt{ar rcs lib\emph{nome}.a}, copiada para dentro do
        executável na ligação;
      \item compartilhada: \texttt{-fPIC -shared}, encontrada na hora de
        executar;
      \item \texttt{-I} headers, \texttt{-L} onde procurar, \texttt{-l} qual
        ligar --- e as bibliotecas por último.
    \end{itemize}
  \end{block}
\end{frame}

% ==================================================================
\section{Os erros que o compilador não pega}
% ==================================================================

\begin{frame}{Quatro erros que atravessam a compilação}
  \begin{center}
  \small
  \begin{tabular}{@{}lll@{}}
    \toprule
    \textbf{Erro} & \textbf{Quando aparece} & \textbf{Sintoma} \\
    \midrule
    \texttt{\#include "matriz.c"}   & na ligação  & \texttt{multiple definition} \\
    Protótipo $\neq$ definição      & ao executar & resultado sem sentido \\
    \texttt{Makefile} sem o header  & ao executar & lixo, ou queda sem motivo \\
    Header sem guarda               & ao compilar & \texttt{redefinition of...} \\
    \bottomrule
  \end{tabular}
  \end{center}
  \begin{block}{}
    \small Os dois do meio são os caros: o programa compila, liga, roda --- e
    mente.
  \end{block}
\end{frame}

\begin{frame}[fragile]{Erro: incluir o \texttt{.c}}
\begin{lstlisting}[basicstyle=\ttfamily\scriptsize, numbers=none]
#include "matriz.c"      /* funciona... ate voce compilar matriz.c tambem */
\end{lstlisting}
  \begin{block}{O que acontece}
    \small O corpo das funções é copiado para dentro do seu \texttt{.c}. Se
    \texttt{matriz.c} também for compilado e ligado, o ligador encontra
    \texttt{cria\_matriz} duas vezes: \texttt{multiple definition}.
  \end{block}
  \begin{block}{Como evitar}
    \small Header nunca tem corpo de função; \texttt{.c} nunca é incluído.
    Se está incluindo um \texttt{.c} para ``funcionar de uma vez'', o que
    faltou foi passar os dois arquivos ao \texttt{gcc}.
  \end{block}
\end{frame}

\begin{frame}[fragile]{Erro: o protótipo que discorda da definição}
  \begin{columns}[T]
    \begin{column}{0.48\textwidth}
      \textbf{\texttt{estat.h}}
\begin{lstlisting}[basicstyle=\ttfamily\tiny, numbers=none]
int media(const double *v, int n);
\end{lstlisting}
    \end{column}
    \begin{column}{0.48\textwidth}
      \textbf{\texttt{estat.c}}
\begin{lstlisting}[basicstyle=\ttfamily\tiny, numbers=none]
double media(const double *v, int n)
{ ... }
\end{lstlisting}
    \end{column}
  \end{columns}
  \begin{block}{Por que ninguém reclama}
    \small \texttt{turma.c} acredita no header e lê o retorno como
    \texttt{int}; \texttt{estat.c} devolve um \texttt{double}. Cada unidade de
    tradução está coerente consigo mesma, e o ligador só compara \emph{nomes}.
  \end{block}
  \begin{alertblock}{A defesa}
    \small \texttt{estat.c} \textbf{inclui} \texttt{estat.h}. Aí o conflito
    passa a acontecer dentro de um arquivo só --- e vira erro de compilação.
  \end{alertblock}
\end{frame}

\begin{frame}[fragile]{Erro: o \texttt{Makefile} que esquece o header}
\begin{lstlisting}[style=makeStyle]
main.o: main.c            # faltou turma.h
\end{lstlisting}
  \begin{block}{O estrago}
    \small Você acrescenta um campo no meio da \texttt{struct Turma}.
    \texttt{turma.o} é refeito com o novo desenho; \texttt{main.o}, não.
    Os dois passam a discordar de onde cada campo começa --- e o
    \texttt{make} diz que está tudo em dia.
  \end{block}
  \begin{block}{As saídas}
    \small \texttt{make clean \&\& make} resolve na hora. Para não depender de
    memória, \texttt{gcc -MMD -MP} gera os arquivos \texttt{.d} com as
    dependências de header, e o \texttt{Makefile} os inclui.
  \end{block}
\end{frame}

\begin{frame}{Recapitulando: armadilhas}
  \begin{block}{O que aprendemos}
    \begin{itemize}
      \item \texttt{.h} declara, \texttt{.c} define, e \texttt{.c} não se
        inclui;
      \item todo \texttt{.c} inclui o próprio \texttt{.h}: é a única forma de
        o compilador conferir o contrato;
      \item dependência de header esquecida no \texttt{Makefile} produz
        bugs que somem quando você recompila do zero --- desconfie deles;
      \item compile sempre com \texttt{-Wall -Wextra}.
    \end{itemize}
  \end{block}
\end{frame}

% ==================================================================
\section{Mãos à obra}
% ==================================================================

\begin{frame}{Altere o programa}
  \begin{block}{Tarefa 1}
    Crie o módulo \texttt{io.h}/\texttt{io.c} com
    \texttt{turma\_le(Turma *t)}, que preenche as notas a partir do teclado, e
    tire essa responsabilidade do \texttt{main.c}. Acrescente as regras
    correspondentes ao \texttt{Makefile}.
  \end{block}
  \begin{block}{Tarefa 2}
    Empacote \texttt{matriz.o} e \texttt{estat.o} em \texttt{libnotas.a} e
    ligue o programa contra ela. Depois apague os \texttt{.o} e recompile só o
    \texttt{main.c}: o executável ainda se forma?
  \end{block}
  \begin{block}{Tarefa 3}
    Acrescente um campo no início da \texttt{struct Turma}, recompile
    \emph{só} \texttt{turma.c} na mão e rode. Explique a saída --- e
    depois conserte pelo \texttt{Makefile}.
  \end{block}
\end{frame}

\begin{frame}{Desafio}
  \begin{block}{Uma biblioteca de verdade, com tipo opaco}
    Reescreva o módulo da turma de modo que \texttt{turma.h} contenha apenas
    \texttt{typedef struct Turma Turma;} --- sem os campos:
    \begin{itemize}
      \item a \texttt{struct} completa fica em \texttt{turma.c}; quem usa a
        biblioteca chega às notas por \texttt{turma\_set} e \texttt{turma\_get};
      \item gere \texttt{libturma.a} e \texttt{libturma.so} dos mesmos
        fontes, com dois alvos no \texttt{Makefile};
      \item escreva um cliente em \emph{outro} diretório, compilado com
        \texttt{-I} e \texttt{-L}, sem copiar nenhum \texttt{.c};
      \item explique por que agora o cliente não consegue depender do
        formato interno da \texttt{struct}.
    \end{itemize}
  \end{block}
\end{frame}

\begin{frame}{Fechamento}
  \begin{columns}[T]
    \begin{column}{0.5\textwidth}
      \textbf{O que mudou hoje}\\[0.3em]
      O programa é o mesmo e imprime o mesmo. O que ganhamos foi mexer em
      um pedaço sem tocar --- nem recompilar --- os outros.
    \end{column}
    \begin{column}{0.5\textwidth}
      \textbf{O fio da aula}\\[0.3em]
      Header é contrato, \texttt{.o} é dívida e oferta, ligador é o
      casamenteiro, \texttt{make} é a memória das datas e biblioteca é o
      pacote pronto.
    \end{column}
  \end{columns}
  \begin{block}{Para levar}
    \begin{itemize}
      \item uma responsabilidade por módulo; o \texttt{.h} é a única porta;
      \item cada \texttt{.c} compila sozinho --- o erro aparece na ligação;
      \item \texttt{-I}, \texttt{-L}, \texttt{-l}: onde estão os headers, onde
        estão as bibliotecas, quais usar;
      \item o bug mais difícil da aula é o que some depois de um
        \texttt{make clean}.
    \end{itemize}
  \end{block}
\end{frame}

\section*{Referências}
\begin{frame}[allowframebreaks]{Referências}
  \nocite{*}
  \bibliographystyle{plain}
  \bibliography{../referencias}
\end{frame}

\end{document}
